\vspace{2mm}
\section{Conclusion}
We introduced \emph{DLO-Lab}, a differentiable simulation platform built on a customized physics solver unifying critical material behaviors, enabling efficient gradient-based policy learning for robotic DLO manipulation. We established a benchmark suite spanning real-world manipulation skills, from precise cable routing to topological unknotting. Addressing the long-horizon nature of these challenges, we developed a specialized DLO agent that employs hierarchical task decomposition and strategic grasp proposals to make complex planning feasible. Our extensive evaluation highlights the trade-offs between gradient-based and sampling-based optimization, offering key insights into policy learning for deformable objects. We validated the physical fidelity of our platform through zero-shot sim-to-real transfer in both open-loop and closed-loop settings: open-loop policies trained via trajectory optimization transfer directly to physical hardware for tasks such as gathering and wiring, while a closed-loop reactive policy achieves $\approx$58\% success on the precision-demanding \textit{Wiring-ring} task. Beyond single-task policies, DLO-Lab can serve as a scalable synthetic data engine: simulation demonstrations collected via trained policies in DLO-Lab enable fine-tuning VLA as a single multi-task policy across tasks, achieving results competitive with per-task RL baselines. We believe DLO-Lab provides a strong foundation for advancing learning-based DLO manipulation, with future opportunities in richer material models and large-scale sim-to-real generalization.